% =====================================================================
% TRADUÇÃO PARA O PORTUGUÊS — versão de trabalho
% Tradução fiel do texto original em inglês, sem alteração de conteúdo,
% argumento ou estrutura. Os marcadores de revisão (\entra, \corr) foram
% removidos e seu conteúdo mantido. Os comentários de revisão estão em
% artigo/revisao_estrutura_coesao.md, separados do texto.
% =====================================================================
\documentclass[conference]{IEEEtran}
\IEEEoverridecommandlockouts

\usepackage[utf8]{inputenc}
\usepackage[T1]{fontenc}
\usepackage[brazil]{babel}
\addto\captionsbrazil{\renewcommand{\figurename}{Fig.}}
\usepackage{cite}
\usepackage{amsmath,amssymb,amsfonts}
\usepackage{graphicx}
\usepackage{tikz}\usetikzlibrary{positioning,arrows.meta,calc,shapes.geometric,fit,backgrounds}
\usepackage{textcomp}
\usepackage{xcolor}
\definecolor{verfill}{RGB}{214,236,216}\definecolor{verdraw}{RGB}{38,110,52}
\usepackage{booktabs}
\usepackage{listings}
\usepackage{courier}
\usepackage{url}
\usepackage[hidelinks]{hyperref}

\setlength{\textfloatsep}{5pt plus 1pt minus 1pt}
\setlength{\floatsep}{4pt plus 1pt minus 1pt}
\setlength{\intextsep}{4pt plus 1pt minus 1pt}
\setlength{\abovedisplayskip}{3pt}
\setlength{\belowdisplayskip}{3pt}
\setlength{\abovecaptionskip}{2pt}
\setlength{\belowcaptionskip}{0pt}
\let\oldbibliography\thebibliography
\renewcommand{\thebibliography}[1]{\oldbibliography{#1}\setlength{\itemsep}{-1pt}\setlength{\parskip}{0pt}\setlength{\parsep}{0pt}}
\usepackage{titlesec}
\usepackage[most]{tcolorbox}
\newtcolorbox{rqbox}{
  colback=black!4,
  colframe=black!55,
  boxrule=0.6pt,
  arc=3pt,
  left=6pt,right=6pt,top=3pt,bottom=3pt,
  boxsep=1pt,
  enlarge top by=1pt, enlarge bottom by=1pt
}

\titleformat{\subparagraph}[runin]{\normalfont\normalsize\bfseries}{\thesubparagraph}{1em}{}
\titlespacing*{\section}{0pt}{4pt}{1.5pt}
\titlespacing*{\subsection}{0pt}{3pt}{1pt}
\titlespacing*{\paragraph}{0pt}{2pt}{6pt}
\linespread{0.96}
\renewcommand{\topfraction}{0.95}
\renewcommand{\bottomfraction}{0.95}
\renewcommand{\textfraction}{0.05}
\renewcommand{\floatpagefraction}{0.85}

\lstset{
  basicstyle=\ttfamily\scriptsize,
  breaklines=true,
  frame=single,
  language=C,
  numbers=left,
  numberstyle=\tiny,
  keywordstyle=\bfseries,
  commentstyle=\itshape,
  literate={á}{{\'a}}1 {é}{{\'e}}1 {í}{{\'i}}1 {ó}{{\'o}}1 {ú}{{\'u}}1 {ã}{{\~a}}1 {õ}{{\~o}}1 {ç}{{\c c}}1
}
\renewcommand{\lstlistingname}{Listagem}

\renewcommand{\baselinestretch}{0.97}

\hypersetup{pdftitle={Verificação Formal de Módulos Críticos do PX4 Autopilot usando Verificação de Modelos Limitada},pdfauthor={Marcela Alves Rodrigues, Lucas C. Cordeiro, Jose Reginaldo H. de Carvalho}}
\begin{document}
\renewcommand{\IEEEkeywordsname}{Palavras-chave}

\title{Verificação Formal de Módulos Críticos do PX4 Autopilot
       usando Verificação de Modelos Limitada\\
}

\author{
\IEEEauthorblockN{Marcela Alves Rodrigues}
\IEEEauthorblockA{\textit{Programa de Pós-Graduação em}\\
\textit{Engenharia Elétrica (PPGEE)} \\
\textit{Universidade Federal do Amazonas}\\
Manaus, Brasil \\
marcela.rodrigues@ufam.edu.br}
\and
\IEEEauthorblockN{Lucas C. Cordeiro}
\IEEEauthorblockA{\textit{Departamento de Eletrônica e}\\
\textit{Computação} \\
\textit{Universidade Federal do Amazonas}\\
Manaus, Brasil \\
lucascordeiro@ufam.edu.br}
\and
\IEEEauthorblockN{Jos\'e Reginaldo H. de Carvalho}
\IEEEauthorblockA{\textit{Instituto de Computação} \\
\textit{Universidade Federal do Amazonas}\\
Manaus, Brasil \\
reginaldo@icomp.ufam.edu.br}
}

\maketitle

\begin{abstract}
O firmware de piloto automático de VANTs em produção pode conter
defeitos latentes que sobrevivem aos testes unitários, com
consequências diretas para a segurança de voo em controladores sem
MMU; a garantia de segurança exige evidência de ausência de defeitos,
e não apenas relatos de bugs. Os esforços anteriores de verificação
formal têm como alvo a lógica de missão ou pilotos automáticos de
pesquisa, deixando o código de drivers em produção amplamente sem
exame. Avaliamos uma metodologia de verificação em oito etapas,
baseada em verificação de modelos limitada com o ESBMC
(\textit{Efficient SMT-based Context-Bounded Model Checker}),
\textit{stubs} de hardware escritos manualmente e asserções derivadas
de \textit{datasheets}, em doze \textit{harnesses} derivados de
trechos de doze módulos críticos para a segurança do PX4 Autopilot
(ramo \texttt{main}, \textit{commit} \texttt{8be2590c15}), com a
revisão da comunidade como parte da metodologia. O resultado principal
é uma garantia limitada, em nível de modelo: os \textit{harnesses} de
dez módulos, que abrangem classes representativas de entradas não
confiáveis, são verificados livres de violações de segurança de
memória e de \textit{overflow} até o limite de desdobramento
(${\approx}94$\,s no total, sem a execução do GPS). O resultado
secundário é a análise de três contraexemplos: um falso positivo no
contador da FIFO (\textit{first-in, first-out}) do BMI088, refutado
por inspeção do código-fonte; uma observação de temperatura do
BMI088, recusada pelos mantenedores; e um \textit{underflow} de
inteiro sem sinal no \textit{harness} do driver uBlox de GPS
(\textit{Global Positioning System}), que se propaga para um
\texttt{memcpy} fora dos limites em um caminho que nenhum teste
unitário do PX4 exercita, proposto ao projeto original; o mantenedor
do PX4 mostrou em seguida que ele é inalcançável no driver de
produção e sugeriu tratar a verificação como um reforço defensivo.
Os \textit{harnesses} expõem estados alcançáveis no modelo que o
código de produção pode nunca alcançar; por isso, os contraexemplos
precisam ser confrontados com o código de produção e com a
comunidade. O ESBMC detecta a violação do GPS com o Boolector (o Z3
não deu veredito dentro do tempo limite); o CBMC (\textit{C Bounded
Model Checker}) a reporta apenas com um \textit{harness} posterior,
semeado, e com um limite de desdobramento reduzido (seu
\textit{harness} original retornou \textsc{verification
successful}); assim, o achado é sensível à configuração da ferramenta
e à pré-condição, o que é uma evidência relevante para a escolha do
conjunto de ferramentas.
\end{abstract}

\begin{IEEEkeywords}
verificação de modelos limitada, ESBMC, PX4, VANT, verificação formal, software embarcado
\end{IEEEkeywords}


\section{Introdução}
\label{sec:introduction}

Veículos Aéreos Não Tripulados (VANTs) empregados em missões críticas
para a segurança exigem um firmware de piloto automático confiável,
cujas falhas têm consequências de segurança bem
documentadas~\cite{Wanner2024}.
A árvore \texttt{src} do PX4 soma ${\approx}710$ mil linhas de código
(LoC) em C/C++ (fontes e cabeçalhos, \texttt{wc\,-l} no \textit{commit}
fixado), executa em microcontroladores sem unidade de gerenciamento de
memória (MMU) sob restrições de tempo real rígido e se comunica com
hardware que pode falhar~\cite{PX4_GitHub,Wanner2024}.
Os testes tradicionais deixam passar falhas disparadas por valores
extremos de sensores ou por casos-limite de aritmética sem
sinal~\cite{Khatiri2024}; também não conseguem demonstrar a ausência
dessas falhas, que é a evidência exigida pela garantia de segurança;
além disso, um verificador que forneça essa evidência precisa
distinguir vulnerabilidades genuínas de artefatos do próprio
\textit{harness} de verificação.

Os esforços existentes de verificação formal de pilotos automáticos
têm como alvo plataformas de pesquisa ou a lógica no nível de
missão~\cite{Pollien2023,Bernardeschi2023}, ou testam o PX4 por meio
de simulação~\cite{Khatiri2024}, enquanto a Verificação de Modelos
Limitada (\textit{Bounded Model Checking}, BMC) se concentra em
código C para sistemas \textit{desktop} ou embarcados~\cite{Sousa2026};
o código de drivers de pilotos automáticos em produção, no entanto,
continua sem exame (Seção~\ref{sec:related}). Uma extensão da BMC
original~\cite{Biere1999} pode fechar essa lacuna: o ESBMC
(\textit{Efficient SMT-based Context-Bounded Model
Checker})~\cite{Cordeiro2012,Monteiro2022} oferece suporte a C e C++,
com base em resolvedores de Satisfatibilidade Módulo Teorias
(SMT)~\cite{deMoura2008,Niemetz2023}. Nosso objetivo é estabelecer
evidência de correção limitada para \textit{harnesses} derivados de
trechos de doze módulos do PX4 Autopilot (ramo \texttt{main},
\textit{commit} \texttt{8be2590c15}) a um custo aceitável, expondo
quaisquer defeitos ao longo do caminho, examinando a latência por
módulo (QP1), as violações que os testes unitários existentes não
detectam (QP2) e as configurações necessárias por módulo (QP3)
(Seção~\ref{sec:results}).

Até onde sabemos, esta é uma das primeiras aplicações de BMC
sistemática ao código do PX4 Autopilot; a metodologia se apoia em
\textit{stubs} de hardware escritos manualmente e em asserções
personalizadas derivadas de \textit{datasheets}
(Seção~\ref{sec:methodology}). Os \textit{harnesses} de dez dos doze
módulos críticos para a segurança são verificados livres de violações
de segurança de memória e de \textit{overflow} até o limite
(${\approx}94$\,s no total, nenhum módulo cronometrado acima de
noventa segundos); a campanha também revelou uma violação alcançável
no modelo e dois não defeitos triados, sendo a violação do GPS
(\textit{Global Positioning System}) reproduzida por outro
verificador, o CBMC (\textit{C Bounded Model Checker}), com um
\textit{harness} posterior, semeado.
Nossas principais contribuições originais são:
(i)~evidência de verificação limitada, em nível de modelo, para
\textit{harnesses} de dez módulos do PX4 que abrangem classes
representativas de entradas não confiáveis;
(ii)~uma metodologia de verificação em oito etapas, com validação
cruzada entre ferramentas e validação de correções;
(iii)~uma violação de segurança de memória alcançável no modelo, no
\textit{harness} do driver uBlox de GPS (\textit{Common Weakness
Enumeration}, CWE-191, CWE-119, Issue~\#26866), proposta ao projeto
original com uma verificação de limites (PR~\#27576) e, em seguida,
mostrada inalcançável no driver de produção pelo mantenedor do PX4,
que enquadrou a verificação como um reforço defensivo;
(iv)~a análise de achados não confirmados: um falso positivo no
contador da FIFO (\textit{first-in, first-out}) do BMI088, refutado
por inspeção do código-fonte (que também havia sido proposto ao
projeto original como uma saturação, no PR~\#27575), e uma observação
de temperatura do BMI088 reportada ao projeto original
(Issue~\#26865) e recusada pelos mantenedores (PR~\#27575, que
continha as duas saturações), esta última reproduzida no ArduPilot
(Issue~\#33272), onde um mantenedor a encerrou;
(v)~um artefato público com todos os \textit{harnesses} e
\textit{logs}
(\url{https://github.com/marcela2881/dissertacao-px4-esbmc}).

\section{Fundamentação Teórica}
\label{sec:background}

\subsection{Verificação de Modelos Limitada}
\label{subsec:bmc}
A BMC~\cite{Biere1999} busca violações de propriedades até a
profundidade~$k$, codificando o programa como uma fórmula de
satisfatibilidade booleana (SAT)/SMT:
\begin{equation}
\mathit{BMC}_\phi(k) =
I(s_0) \wedge
\bigwedge_{i=0}^{k-1} T(s_i, s_{i+1}) \wedge
\bigvee_{i=0}^{k} \neg\phi(s_i).
\label{eq:bmc}
\end{equation}
Uma atribuição que satisfaz a fórmula é um contraexemplo concreto. A
BMC é \emph{correta} e \emph{completa} até o limite $k$: as violações
reportadas são reais \emph{para o modelo codificado}; a ausência de
violações até~$k$ não é uma prova.

\subsection{O Verificador ESBMC}
\label{subsec:esbmc}

O ESBMC~\cite{Cordeiro2012,Monteiro2022,Menezes2024} é um verificador
BMC de código aberto para C e C++.
Seu \textit{frontend} baseado no Clang traduz o código-fonte para um
programa GOTO; a execução simbólica desdobra os laços até o
limite~$k$ e emite condições de verificação (VCCs), resolvidas pelo
Z3~\cite{deMoura2008}, pelo Boolector~\cite{Brummayer2009} ou pelo
Bitwuzla~\cite{Niemetz2023}.
As propriedades definidas pelo usuário são expressas como asserções
no próprio código.

\subsection{O PX4 Autopilot}
\label{subsec:px4}

O PX4~\cite{PX4_GitHub,PX4_Documentation} é um controlador de voo de
código aberto para aeronaves multirrotor, de asa fixa e de decolagem
e pouso verticais (VTOL).
Ele executa sobre o NuttX em microcontroladores STM32 sem MMU, nos
quais todos os módulos compartilham um único espaço de endereçamento;
assim, um defeito de memória em qualquer driver pode corromper
estados críticos para o voo~\cite{PX4_Documentation}; esses defeitos
também são uma preocupação de segurança da
informação~\cite{Siddiqi2022}.
Os módulos se comunicam por meio do \textit{middleware} de
publicação--assinatura uORB (\textit{micro Object Request Broker})
(cf. Fig.~\ref{fig:px4_arch}).

\begin{figure}[t]
    \centering
    \resizebox{0.88\columnwidth}{!}{%
\begin{tikzpicture}[font=\fontsize{7}{8}\selectfont,
 mod/.style={draw,rounded corners=1.5pt,minimum height=6mm,text width=27mm,align=center,inner sep=1pt,fill=white,draw=black!55},
 ver/.style={mod,fill=verfill,draw=verdraw,line width=0.8pt,text=black},
 bus/.style={draw,rounded corners=2pt,minimum height=19mm,minimum width=17mm,align=center,fill=white,thick,inner sep=1pt},
 grp/.style={draw=black!45,dashed,rounded corners=2pt,inner sep=2mm},
 ar/.style={-{Stealth[length=1.5mm]},line width=0.4pt}]
\node[ver] (gps) {Driver de GPS (u-blox)};
\node[ver,above=3.4mm of gps] (imu) {Driver da IMU (BMI088)};
\node[ver,below=3.4mm of gps] (rc) {Entrada de RC / módulo Expo};
\node[bus,right=7mm of gps] (uorb) {\textbf{uORB}\\[-1pt]{\fontsize{6}{7}\selectfont barramento}\\[-1pt]{\fontsize{6}{7}\selectfont de mensagens}};
\node[mod,right=7mm of uorb] (fc) {Controle de voo};
\node[mod,above=3.4mm of fc] (ekf) {EKF / estimador de estado};
\node[ver,below=3.4mm of fc] (mav) {Módulo MAVLink};
\node[above=1.3mm of imu,align=center] (pl) {\textbf{Publicadores}\\[-1.5pt]{\fontsize{5.5}{6}\selectfont\texttt{orb\_publish()}}};
\node[above=1.3mm of ekf,align=center] (sl) {\textbf{Assinantes}\\[-1.5pt]{\fontsize{5.5}{6}\selectfont\texttt{orb\_subscribe()}}};
\begin{scope}[on background layer]
\node[grp,fit=(imu)(gps)(rc)(pl)] (gp) {};
\node[grp,fit=(ekf)(fc)(mav)(sl)] (gs) {};
\end{scope}
\draw[ar](imu.east)--(uorb.north west);
\draw[ar](gps.east)--(uorb.west);
\draw[ar](rc.east)--(uorb.south west);
\draw[ar](uorb.north east)--(ekf.west);
\draw[ar](uorb.east)--(fc.west);
\draw[ar](uorb.south east)--(mav.west);
\node[ver,minimum height=3.6mm,minimum width=6mm,text width=0pt,inner sep=0pt,anchor=north west] (lg) at ([yshift=-2.5mm]gp.south west) {};
\node[right=1.2mm of lg,font=\fontsize{6.5}{7.5}\selectfont,anchor=west] {módulo coberto por um \textit{harness}};
\end{tikzpicture}
}
    \caption{Arquitetura de comunicação do PX4: os drivers publicam
    dados de sensores, por meio do uORB, para os módulos de
    estimação, controle e comunicação~\cite{PX4_Documentation}. Os
    módulos destacados em verde são um subconjunto dos doze módulos
    cobertos neste trabalho, exibidos assim por legibilidade.}
    \label{fig:px4_arch}
\end{figure}

\section{Trabalhos Relacionados}
\label{sec:related}

\paragraph{BMC de firmware e de código C/C++}
Ferramentas de BMC como o ESBMC e o CBMC são comumente avaliadas em
\textit{benchmarks} de competição curados~\cite{Beyer2022}.
Sousa et al.~\cite{Sousa2026} aplicaram o ESBMC a grandes projetos C
de código aberto; herdamos o fluxo de BMC desse trabalho, mas
acrescentamos duas etapas específicas para firmware: a abstração da
E/S de hardware baseada em \textit{stubs} (UART, \textit{Universal
Asynchronous Receiver-Transmitter}; I2C, \textit{Inter-Integrated
Circuit}; SPI, \textit{Serial Peripheral Interface}; uORB) e as
asserções personalizadas derivadas de \textit{datasheets}.

\paragraph{Verificação e segurança de VANTs}
A verificação formal de firmware de pilotos automáticos continua
sendo uma lacuna em aberto~\cite{Luckcuck2019}.
Pollien~\cite{Pollien2023} verificou o Paparazzi com Frama-C e Coq;
Khatiri et al.~\cite{Khatiri2024} propuseram testes do PX4 baseados
em simulação, complementares à BMC.
Bernardeschi et al.~\cite{Bernardeschi2023} aplicaram métodos formais
no nível da lógica de missão, e não no código de drivers.
Diferentemente de todos esses trabalhos, nosso alvo é o código do
PX4 Autopilot em um \textit{commit} fixado do ramo \texttt{main},
analisado, por meio de \textit{harnesses}, quanto à ausência limitada
de violações de segurança de memória e de \textit{overflow}, em vez
de testado ou monitorado em tempo de execução.

\section{Metodologia}
\label{sec:methodology}

Propomos uma metodologia em oito etapas (Fig.~\ref{fig:methodology}):
(i)~seleção de módulos, (ii)~\textit{stubs} de abstração de hardware,
(iii)~especificação de propriedades de segurança, (iv)~execução do
ESBMC, (v)~análise de contraexemplos, (vi)~relato à comunidade,
(vii)~validação cruzada entre ferramentas e verificação de correções,
e (viii)~portabilidade entre pilotos automáticos e integração
contínua/entrega contínua (CI/CD).
Adotamos uma estratégia modular porque a árvore \texttt{src} do PX4
(${\approx}710$ mil~LoC) torna a verificação monolítica impraticável.
As etapas~(i) e (iv)--(vi) seguem o fluxo de trabalho de
Sousa et al.~\cite{Sousa2026}; nossas adições são as etapas
específicas para firmware~(ii)--(iii) (\textit{stubs} de abstração de
hardware e asserções derivadas de \textit{datasheets}) e as
etapas~(vii)--(viii), que acrescentam validação cruzada entre
ferramentas, reverificação de correções, verificação entre pilotos
automáticos e integração com CI/CD.
Na prática: a etapa (v) confronta cada contraexemplo com o código-fonte
não modificado do driver, por reexecução ou por leitura do código
(Seção~\ref{subsec:case-gps}); a etapa (vi) reporta os achados à
comunidade do PX4 (no GPS e na FIFO, a \textit{issue} aberta no
projeto original pediu a triagem, e a reexecução veio depois,
Seção~\ref{subsec:case-imu}); a etapa (vii) faz a verificação cruzada
com um verificador independente e reverifica qualquer correção
(\S\ref{subsec:overall}); a etapa (viii) verifica a reprodução em
outro piloto automático e executa os \textit{harnesses} na CI a cada
\textit{push} (\S\ref{subsec:discussion}).

\begin{figure}[t]
\centering
\resizebox{0.9\columnwidth}{!}{%
\begin{tikzpicture}[font=\fontsize{7}{8}\selectfont, node distance=1.4mm,
 st/.style={draw,rounded corners=1.2pt,fill=white,minimum height=5mm,text width=37mm,align=center,inner sep=1pt},
 add/.style={st,fill=verfill,draw=verdraw,line width=0.8pt},
 outb/.style={draw,rounded corners=1.2pt,minimum height=5mm,text width=19mm,align=center,inner sep=1pt},
 ar/.style={-{Stealth[length=1.6mm]},line width=0.4pt}]
\node[st] (s1) {1. Seleção dos módulos críticos};
\node[st,below=of s1] (s2) {2. Instrumentação com \textit{stubs}};
\node[st,below=of s2] (s3) {3. Especificação das propriedades};
\node[st,below=of s3] (s4) {4. Configuração e execução do ESBMC};
\node[st,below=of s4] (s5) {5. Análise de contraexemplos};
\node[st,right=10mm of s1] (s6) {6. Relato à comunidade do PX4};
\node[add,below=of s6] (s7) {7. Validação cruzada (CBMC) + verificação da correção};
\node[add,below=of s7] (s8) {8. Outro piloto automático (ArduPilot) + CI/CD};
\node[diamond,draw,aspect=2.4,inner sep=0pt,minimum width=14mm,below=4mm of s5,fill=white] (d) {Resultado?};
\node[outb,fill=black!5,draw=black!55] (o1) at ($(d)+(-27mm,-12mm)$) {Bug confirmado};
\node[outb,fill=black!5,draw=black!55] (o2) at ($(d)+(0,-12mm)$) {Falso positivo};
\node[outb,fill=black!5,draw=black!55] (o3) at ($(d)+(27mm,-12mm)$) {Verificado OK};
\foreach \a/\b in {s1/s2,s2/s3,s3/s4,s4/s5,s5/d,s6/s7,s7/s8} \draw[ar](\a)--(\b);
\draw[ar](d.west)-|(o1.north); \draw[ar](d.south)--(o2.north); \draw[ar](d.east)-|(o3.north);
\node[font=\fontsize{6}{7}\selectfont,anchor=south east] at ([xshift=-1mm]d.west) {bug real};
\node[font=\fontsize{6}{7}\selectfont,anchor=south west] at ([xshift=1mm]d.east) {verificado};
\coordinate (g) at ($(s5.east)+(5mm,0)$);
\draw[line width=0.4pt](o1.south)--++(0,-2.5mm)-|(o3.south);
\draw[line width=0.4pt](o3.east)--++(2.5mm,0)|-(g);
\draw[ar](g)|-(s6.west);
\end{tikzpicture}
}
\caption{Fluxo de verificação em oito etapas. As etapas~1 e 4--6
    seguem o fluxo de trabalho de Sousa et al.~\cite{Sousa2026}; as
    etapas~2--3 (\textit{stubs} de hardware e asserções derivadas de
    \textit{datasheets}) e as etapas sombreadas~7--8 são as adições
    deste trabalho. Nenhum dos três contraexemplos analisados é
    classificado como bug confirmado
    (Seções~\ref{subsec:case-imu} e~\ref{subsec:case-gps}).}
    \label{fig:methodology}
\end{figure}

Como exemplo, o driver uBlox de GPS instancia todas as etapas: ele é
selecionado pelo impacto na segurança (i), sua E/S é substituída por
\textit{stubs} (ii), a propriedade de limites do \textit{buffer} é
especificada no \texttt{memcpy} (iii), o ESBMC gera um contraexemplo
(iv--v), que é registrado no projeto original (vi) como
Issue~\#26866; a revisão do mantenedor e a nossa análise do
código-fonte da lógica de verificação de limites do driver mostram,
então, que o estado violador não ocorre em produção; o achado é
verificado de forma cruzada com o CBMC e reverificado na correção
(vii) e registrado para a CI (viii); a Seção~\ref{subsec:case-gps}
detalha esse rastro e sua violação.

\subsection{Seleção de Módulos e Construção dos \textit{Harnesses}}
\label{subsec:harnesses}
Os módulos são priorizados pelo impacto na segurança ao longo do
caminho do sensoriamento à atuação e pela viabilidade de isolamento
com \textit{stubs}; selecionar pela viabilidade dos \textit{stubs}
introduz um viés de seleção (Seção~\ref{sec:threats}).
Escrevemos manualmente \textit{harnesses} para doze módulos, que
totalizam ${\approx}2.800$ linhas (\texttt{wc\,-l} sobre os arquivos
\texttt{test\_*.cpp}), resumidos por módulo na
Tabela~\ref{tab:benchmarks}.

\begin{table}[t]
\centering
\caption{Módulos cobertos: tamanho do alvo, tamanho do
\textit{harness}, localização no código-fonte do PX4 e justificativa
da seleção (classe de entrada exercitada).}
\label{tab:benchmarks}
\scriptsize
\setlength{\tabcolsep}{3pt}
\resizebox{\columnwidth}{!}{%
\begin{tabular}{@{}llrrl@{}}
\toprule
\textbf{Módulo} & \textbf{Caminho no PX4 (\texttt{src/})} & \textbf{Alvo} & \textbf{Harn.} & \textbf{Classe de entrada}\\
\midrule
GPS uBlox      & \texttt{drivers/gps}          & 1712 & 294 & Sensor serial \\
IMU BMI088     & \texttt{drivers/imu/.../bmi088} & 700 & 319 & Sensor serial \\
MAVLink        & \texttt{modules/mavlink}       & 3624 & 82 & Enlace de telemetria \\
Expo RC        & \texttt{lib/mathlib/math}      & 303 & 129 & Entrada de rádio \\
Monitor de bateria & \texttt{lib/battery}       & 449 & 242 & Sensor de energia \\
Saída PWM      & \texttt{drivers/pwm\_out}      & 337 & 208 & Atuação \\
RC Update      & \texttt{lib/rc}                & 769 & 253 & Entrada de rádio \\
Navigator      & \texttt{modules/navigator}     & 497 & 291 & Lógica de missão \\
Cov. do EKF2   & \texttt{modules/ekf2/EKF}      & 385 & 296 & Estimação de estado \\
Magnetômetro   & \texttt{modules/sensors}       & 792 & 239 & Sensor serial \\
Barômetro      & \texttt{modules/sensors}       & 246 & 278 & Sensor serial \\
VTOL           & \texttt{modules/vtol\_att\_control} & 591 & 146 & Atuação \\
\bottomrule
\multicolumn{5}{l}{Alvo = LoC dos arquivos-fonte do PX4 dos quais o \textit{harness} extrai trechos}\\
\multicolumn{5}{l}{(\textit{commit} fixado); Harn. = LoC do \textit{harness} (ambos com \texttt{wc\,-l});}\\
\multicolumn{5}{l}{RC = radiocontrole; IMU = unidade de medição inercial;}\\
\multicolumn{5}{l}{EKF2 = Filtro de Kalman Estendido 2; PWM = modulação por largura de pulso.}\\
\multicolumn{5}{l}{Módulos selecionados pelo impacto na segurança e pela viabilidade dos \textit{stubs}.}
\end{tabular}}
\end{table}
A instrumentação de um módulo exige três artefatos escritos
manualmente: um \textit{harness} de ponto de entrada, que aciona o
módulo com entradas não determinísticas; \textit{stubs} que
substituem a E/S de hardware e de \textit{middleware} (UART, I2C,
SPI, uORB); e as asserções derivadas de \textit{datasheets}. Os três
ficam no mesmo arquivo \texttt{test\_*.cpp} de cada módulo; o
repositório não tem arquivos de \textit{stubs} separados. Em relação
a ${\approx}10.400$ linhas de código-alvo, isso corresponde a uma
razão próxima de $1{:}3{,}7$, com os valores por módulo na
Tabela~\ref{tab:benchmarks}. Os fontes do PX4 não foram modificados.

Em vez de um percentual de cobertura de linhas, que induziria a erro
quando cada \textit{harness} substitui por \textit{stubs} a maior
parte do seu módulo e explora apenas os caminhos alcançáveis com
$k=10$, caracterizamos a cobertura por \emph{categoria de entrada}
(Seção~\ref{subsec:overall}). A E/S de hardware, as chamadas ao NuttX
e as dependências entre módulos são substituídas por \textit{stubs}
que retornam valores \texttt{nondet\_T()}.
Cada \textit{harness} fornece um ponto de entrada com \textit{stubs}
e asserções.
O \textit{harness} do uBlox de GPS, por exemplo, é uma das cinco
funções de teste do seu arquivo (\texttt{test\_choice=3} em
\texttt{main}) e é a função que produz o contraexemplo analisado na
Seção~\ref{subsec:case-gps}: ele sorteia valores não determinísticos
para \texttt{input\_len} e \texttt{dump\_buffer.len}, restringe-os
com \texttt{\_\_ESBMC\_assume} a
\texttt{input\_len}$\,\in\,(0,\text{GPS\_DUMP\_DATA\_SIZE}{+}10]$ e
\texttt{dump\_buffer.len}$\,\geq\,\text{GPS\_DUMP\_DATA\_SIZE}{-}5$,
chama \texttt{dumpGpsData()} e afirma que o comprimento resultante do
\textit{buffer} permanece dentro de \texttt{GPS\_DUMP\_DATA\_SIZE}
(essa asserção final é uma tautologia, pois seu último disjunto
$(\texttt{len}\,\&\,\texttt{0x7F})\le 200$ é sempre verdadeiro;
assim, a violação é capturada pela verificação automática de
desreferência do ESBMC dentro do \texttt{memcpy}). Esses limites
admitem os valores do contraexemplo (\texttt{input\_len=10},
\texttt{dump\_buffer.len=207}), que levam ao \texttt{memcpy} fora dos
limites mostrado na Listagem~\ref{lst:bug-gps}.
Os fontes dos \textit{harnesses} estão no repositório de
artefatos\footnote{\url{https://github.com/marcela2881/dissertacao-px4-esbmc}},
que contém os doze \textit{harnesses} de teste (um arquivo
\texttt{test\_*.cpp} por módulo), os \textit{logs} brutos do ESBMC de
todas as execuções aqui reportadas, os \textit{scripts} e
\textit{logs} da validação cruzada com o CBMC, as linhas de comando
por módulo e um guia de reprodução; o repositório ainda não contém
arquivos separados de \textit{drivers} de reexecução nem
\textit{logs} de \textit{sanitizers} (a reexecução é descrita em
texto, Seção~\ref{subsec:case-gps}).

\subsection{Propriedades de Segurança}
\label{subsec:properties}

As propriedades \emph{automáticas} (limites de vetores, segurança de
ponteiros, \textit{overflow} aritmético e divisão por
zero~\cite{Menezes2024,Monteiro2022}) são habilitadas uniformemente.
As propriedades \emph{personalizadas} codificam requisitos dos
\textit{datasheets}.
O \textit{datasheet} do BMI088~\cite{Bosch2022} restringe a
temperatura de operação a $[-40,\,+85]\,^\circ$C. Como a conversão
afim do driver não é saturada, qualquer leitura bruta fora dessa
faixa viola $\phi_{\text{temp}}$ por construção; portanto, tratamos
$\phi_{\text{temp}}$ não como uma garantia de hardware, mas como uma
\emph{sonda de cobertura} que sinaliza \emph{a falta de validação da
faixa de entrada} no driver:
\begin{equation}
  \phi_{\text{temp}}: \quad -40{,}0 \leq \textit{temperatura} \leq 85{,}0
  \label{eq:prop_temp}
\end{equation}
O driver de GPS exige que o nível de preenchimento do \textit{buffer}
de \textit{dump} permaneça dentro da capacidade do registro,
\texttt{sizeof(dump\_data->data)}; o \textit{harness} o modela com a
macro \texttt{GPS\_DUMP\_DATA\_SIZE}:
\begin{equation}
  \phi_{\text{gps}}: \quad \textit{dump\_data.len} \leq \texttt{GPS\_DUMP\_DATA\_SIZE}
  \label{eq:prop_gps}
\end{equation}
Como \texttt{len} é uma entrada livre, verificar $\phi_{\text{gps}}$
na entrada seria trivialmente violado pela própria entrada;
verificamos o acesso no ponto de uso do \texttt{memcpy}, onde o
driver não deve acessar além do \textit{buffer} para nenhum valor de
\texttt{len}.
Os contraexemplos são confrontados com o código-fonte não modificado
para descartar artefatos dos \textit{stubs} (por reexecução ou por
leitura do código); os \textit{drivers} de reexecução ainda não estão
no artefato.

\section{Avaliação Experimental}
\label{sec:results}

Esta seção avalia a metodologia nos \textit{harnesses} de doze
módulos do PX4. Abordamos três questões de pesquisa:
\textbf{QP1}: Qual é a latência de verificação por módulo em hardware
convencional? \textbf{QP2}: A BMC expõe violações alcançáveis no
modelo que os testes unitários existentes do PX4 não exercitam?
\textbf{QP3}: Quais configurações por módulo são necessárias para
tornar cada módulo tratável?

\subsection{Configuração Experimental}
\label{subsec:setup}

A verificação foi executada em um Intel~i7-13700K / 64~GB /
Ubuntu~24.04, sobre os fontes do ramo \texttt{main} do PX4
(\textit{commit}~\texttt{8be2590c15}, 2026-06-04); o artefato
completo (\textit{harnesses} e \textit{logs}) está arquivado no
Zenodo, DOI~\texttt{10.5281/zenodo.23282236}. Os tempos reportados
são a mediana de três execuções. O \textit{commit} está 2.935
\textit{commits} após a v1.16.0 e é posterior ao lançamento da
v1.17.0; os \textit{harnesses} são cópias de trechos do código-fonte,
e não são compilados contra os arquivos do PX4, e algumas constantes
diferem (\textit{buffer} do GPS: 200 bytes no \textit{harness}, 79 no
PX4); assim, as VCCs cobrem os \textit{harnesses}, e não os módulos
completos. Todos os \textit{harnesses} foram verificados com o
ESBMC~v8.1.0, usando o Boolector~3.2.4~\cite{Brummayer2009} como
resolvedor.\footnote{A \textit{string} de versão
\texttt{ESBMC version 8.1.0} foi lida em todos os doze \textit{logs}
de execução; os contraexemplos do BMI088 e do GPS foram registrados
originalmente no projeto original (Issues~\#26865 e \#26866) a partir
de execuções anteriores (ESBMC~v7.10; GPS em 12,1\,s), e os tempos e
valores dos contraexemplos aqui apresentados foram extraídos dos
\textit{logs} da v8.1.0.} O Boolector também é a configuração com a
qual o relato do GPS no projeto original (Issue~\#26866) foi
registrado; em execuções exploratórias anteriores, o Z3 não retornou
veredito no \textit{harness} do GPS (Seção~\ref{sec:threats}).

\noindent\begin{minipage}{\columnwidth}
Invocação padrão:
\begin{lstlisting}[language=bash,breaklines=true]
esbmc module.cpp --unwind 10 --no-unwinding-assertions \
  --overflow-check --default-solver boolector \
  --timeout 120s
\end{lstlisting}
\end{minipage}
Como \texttt{--no-unwinding-assertions} está habilitado, os laços que
excedem $k=10$ são truncados silenciosamente
(Seção~\ref{sec:threats}).
Ajustes: o MAVLink usou \texttt{--unwind~300}; o Navigator/Mission,
\texttt{--unwind~15}; o VTOL, \texttt{--no-library}, que desabilita a
biblioteca C abstrata do ESBMC: \texttt{sinf}/\texttt{sqrtf} passam a
retornar valores não determinísticos, uma sobreaproximação sob a qual
nenhuma violação das demais verificações deixa de ser detectada; as
verificações internas da biblioteca, no entanto, deixam de ser
realizadas, de modo que o veredito cobre proporcionalmente menos
propriedades.

\subsection{Resultados Gerais}
\label{subsec:overall}

A Tabela~\ref{tab:results} resume os resultados. Os
\textit{harnesses} de dez dos doze módulos são verificados sem
violações até o limite; essa evidência limitada, em nível de modelo,
é o resultado principal da campanha (com o escopo qualificado na
Seção~\ref{sec:threats}).
O \textit{harness} do uBlox de GPS gerou a única violação
(Issue~\#26866); o candidato do contador da FIFO do BMI088 foi um
falso positivo na triagem; a observação sobre a faixa de temperatura
continua contestada (Seção~\ref{subsec:case-imu}).
\begin{table}[t]
\centering
\caption{Resultados da verificação dos \textit{harnesses} de doze
módulos do PX4.}
\label{tab:results}
\scriptsize
\setlength{\tabcolsep}{3pt}
\begin{tabular}{@{}lcrccr@{}}
\toprule
\textbf{Módulo} & \textbf{Resultado} & \textbf{VAM}$^\P$ & \textbf{Desd.} & \textbf{Resolv.} & \textbf{Tempo}$^\|$\\
\midrule
IMU BMI088        & \textit{viol.} & 0$^\ddagger$ & 10 & Bool.$^\star$ & 0,1\,s$^{\#}$ \\
GPS uBlox         & \textit{viol.} & 1 & 10 & Bool.$^\star$ & 64,9\,s \\
Parser MAVLink    & \textit{sem viol.} & 0 & 300 & Bool. & 49,2\,s \\
Expo RC           & \textit{sem viol.} & 0 & 10 & Bool. & 4,9\,s \\
Monitor de bateria & \textit{sem viol.} & 0 & 10 & Bool. &   6,7\,s \\
Saída PWM         & \textit{sem viol.} & 0 & 10 & Bool. &   0,2\,s \\
RC Update         & \textit{sem viol.} & 0 & 10 & Bool. &  9,5\,s \\
Navigator/Mission & \textit{sem viol.} & 0 & 15 & Bool. &   0,1\,s \\
Covariância EKF2  & \textit{sem viol.} & 0 & 10 & Bool. & 12,2\,s \\
Magnetômetro      & \textit{sem viol.} & 0 & 10 & Bool. &  6,6\,s \\
Barômetro         & \textit{sem viol.} & 0 & 10 & Bool. &   0,4\,s \\
VTOL$^\S$         & \textit{sem viol.} & 0 & 10 & Bool. &  4,0\,s \\
\midrule
\textbf{Total} & \textit{10 sem viol., 2 viol.} & 1 & --- & --- & ${\approx}94$\,s$^{\S\S}$ \\
\bottomrule
\multicolumn{6}{l}{\textit{sem viol.} = nenhuma violação até $k$; \textit{viol.} = violação encontrada.}\\
\multicolumn{6}{l}{$^\P$VAM = violações alcançáveis no modelo (não são defeitos confirmados).}\\
\multicolumn{6}{l}{$^\star$ESBMC v8.1.0; Boolector~3.2.4 em todos os módulos.}\\
\multicolumn{6}{l}{$^\ddagger$Falso positivo da FIFO; temperatura recusada no projeto original.}\\
\multicolumn{6}{l}{$^{\S\S}$93,8\,s para os dez módulos concluídos; 158,7\,s com o GPS (64,9\,s).}\\
\multicolumn{6}{l}{$^\S$VTOL com \texttt{--no-library}: verificações internas da biblioteca não realizadas.}\\
\multicolumn{6}{l}{$^\|$Mediana de três execuções.}\\
\multicolumn{6}{l}{$^{\#}$Módulos com violação param na primeira violação (tempo de decisão).}
\end{tabular}
\end{table}
Os tempos variaram de $0{,}1$ a $64{,}9$\,s; os tempos de relógio na
CI não foram medidos (Seção~\ref{subsec:discussion}).
\begin{rqbox}
\textbf{Esses resultados respondem à QP1}: todo \textit{harness}
cronometrado é verificado em menos de noventa segundos em hardware
convencional. Como cada um é verificado isoladamente, por meio de um
\textit{harness} pequeno, isso demonstra a \emph{viabilidade} por
módulo, e não a escalabilidade para o sistema inteiro.
\end{rqbox}

Os módulos cobertos não são uma amostra uniforme: eles estão no
caminho crítico para a segurança que vai do sensoriamento à atuação e
cobrem as classes representativas de entradas não confiáveis, a
saber, barramentos seriais de sensores, o enlace de rádio/telemetria,
a atuação, a estimação e a lógica de missão
(Tabela~\ref{tab:benchmarks}).
O restante do código é dominado por código de estação de solo, de
registro (\textit{logging}) e de compilação, com baixa consequência
para a segurança em tempo de execução; as exceções críticas para o
voo (\textit{commander}, \textit{failsafe}, \textit{geofence}) são
trabalhos futuros prioritários (Seção~\ref{sec:conclusion}; lacunas
restantes na Seção~\ref{sec:threats}).

Para avaliar a dependência em relação à ferramenta, executamos o
CBMC~5.95.1~\cite{Kroening2014} com opções equivalentes
(\texttt{--unwind~10}, verificações de \textit{overflow} e de
limites, sem asserções de desdobramento) e o mesmo tempo limite de
$120$\,s, no mesmo \textit{harness} para a IMU~BMI088 (as duas
ferramentas concordaram) e em um \textit{harness} separado para o
uBlox de GPS. Para o uBlox de GPS, com o \textit{harness} posterior
do CBMC (adicionado em 2026-09-24; o original, com
\texttt{len}~$<$~$200$, retornou \textsc{verification successful} na
Issue~\#26866, pois sua pré-condição exclui os estados violadores) e
com a pré-condição \texttt{len}~$\geq$~$190$
(\texttt{GPS\_DUMP\_DATA\_SIZE}$\,-\,$10, semeando o campo de
comprimento perto do limite do \textit{buffer}; o \textit{harness} do
ESBMC usa $-5$), as duas ferramentas sinalizam uma violação, mas o
CBMC apenas com limites reduzidos (\texttt{--unwind~2}, \texttt{3}
e~\texttt{5}: destino do \texttt{memcpy} não gravável); com
\texttt{--unwind~10}, ele não terminou dentro de um limite maior, de
$30$\,min.
O ESBMC reporta um acesso fora dos limites dentro do \texttt{memcpy}
(uma falha de desreferência); classificamos o defeito subjacente como
CWE-191 (o \textit{underflow} sem sinal) levando a CWE-119/787 (a
escrita fora dos limites), uma classificação nossa, e não uma saída
da ferramenta (não é um defeito confirmado,
Seção~\ref{subsec:case-gps}).
A concordância sob a mesma pré-condição corrobora apenas a
codificação, e não que o estado ocorra no driver de produção
(Seção~\ref{subsec:case-gps}).

\subsection{Estudo de Caso: Driver da IMU BMI088}
\label{subsec:case-imu}

O driver do BMI088 (700~LoC) produziu dois candidatos; a análise
levou a conclusões diferentes para cada um.

\subsubsection{Temperatura fora da faixa de operação (contestada, rejeitada no projeto original)}
O driver não valida a temperatura antes da aritmética de compensação
térmica; o \textit{datasheet} do BMI088~\cite{Bosch2022} restringe a
operação a $[-40,\,+85]\,^\circ$C (Eq.~\eqref{eq:prop_temp}).
O ESBMC produziu \texttt{temp\_msb=63}, \texttt{temp\_lsb=95}
(\texttt{Temp\_uint11=506}), resultando em $86{,}25\,^\circ$C
(Listagem~\ref{lst:bug-temp}).
Isso é CWE-1284 (validação inadequada de quantidade): o driver
consome, sem verificação, uma leitura fora da faixa certificada.
O \textit{harness} não precisa de \texttt{\_\_ESBMC\_assume} aqui:
suas duas entradas de 8 bits totalmente não determinísticas já
abrangem a faixa de valores brutos de 11 bits do sensor; assim, a
asserção é exercitada sobre o domínio alcançável do registrador, e
não de forma vazia.
Se isso é um defeito é algo contestado. Na revisão da correção
enviada ao projeto original (PX4~PR~\#27575), um revisor se opôs a
saturar o valor (``falsificar o número'') e perguntou o que o próprio
sensor reporta fora da sua faixa de operação; um mantenedor então
encerrou o PR, porque ele não segue o padrão estabelecido em outros
drivers. A Issue~\#26865 foi encerrada automaticamente como não
planejada.
Mantemos a observação apenas como uma nota de projeto, \emph{sem
consequência de segurança confirmada e sem correção aceita}, que
ilustra a etapa de triagem da metodologia, e não um defeito.

\begin{lstlisting}[float=t,basicstyle=\ttfamily\scriptsize,
numbers=none,language={},captionpos=b,
caption={Contraexemplo do ESBMC: temperatura do BMI088 de $86{,}25\,^\circ$C
$>$ $85\,^\circ$C. Falta de validação da faixa de entrada (CWE-1284).},
label=lst:bug-temp]
[Counterexample]
State 2 -- test_temperature_calculation, line 120
  temp_msb = 63, temp_lsb = 95
State 7 -- updateTemperature, line 42
  Temp_uint11 = 506
State 11 -- updateTemperature, line 51
  temperature = 86.250000f
Violated property: assertion
  temperature >= -40.0f && temperature <= 85.0f

VERIFICATION FAILED
\end{lstlisting}

\subsubsection{Candidato do contador da FIFO (falso positivo)}
O ESBMC produziu \texttt{count=9088}~bytes (mais de $8\times$ a FIFO
de $1.024$~bytes), sugerindo CWE-190/CWE-787.
A confrontação desse contraexemplo com o driver de produção
\emph{não instrumentado} mostrou um \textbf{falso positivo}. O driver
lançado protege a leitura em três níveis: o contador de bytes é
comparado com o tamanho da FIFO e reinicializado sempre que
\texttt{fifo\_byte\_counter}$\geq$\texttt{FIFO::SIZE}; o número de
amostras é limitado a \texttt{FIFO\_MAX\_SAMPLES}; e o comprimento da
transferência é saturado por
\texttt{math::min(samples*sizeof(FIFO::DATA)\allowbreak+4,~FIFO::SIZE)}.
O valor $9088\geq1.024$ dispara a reinicialização da primeira
proteção; portanto, nenhum acesso fora dos limites é alcançável. O
candidato surgiu porque o \textit{harness} omitiu essas proteções.
Uma saturação para ele também foi proposta ao projeto original no
mesmo PR~\#27575, mas a inspeção do driver de produção mostra que o
candidato é um falso positivo. A confrontação com o código-fonte de
produção o detectou.

\subsection{Estudo de Caso: Driver uBlox de GPS}
\label{subsec:case-gps}

O \textit{harness} do uBlox (que extrai trechos de um arquivo-fonte
de $1.712$~LoC) foi verificado em $64{,}9$\,s.
A verificação de espaço disponível subtrai o nível de preenchimento
do \textit{buffer} da sua capacidade (o \textit{harness} define
\texttt{GPS\_DUMP\_DATA\_SIZE} como $200$; o driver lançado declara
\texttt{uint8[79]~data}):
\begin{lstlisting}[language=C]
if (write_len > GPS_DUMP_DATA_SIZE - dump_data->len)
\end{lstlisting}
Quando \texttt{dump\_data->len} excede a capacidade, a diferença é
negativa ($-7$ para $\texttt{len}=207$: o comprimento
\texttt{uint8\_t} é promovido a \texttt{int}, de modo que a subtração
é com sinal no \textit{harness}) e é convertida em um valor
\texttt{size\_t} enorme na comparação com \texttt{write\_len}, o que
torna a proteção sempre falsa.
O ESBMC produziu \texttt{input\_len=10} e
\texttt{dump\_data->len=207}; a proteção é então sempre falsa, e o
\texttt{memcpy} seguinte escreve fora do \textit{buffer}
(Listagem~\ref{lst:bug-gps}).
CWE-191 (\textit{underflow} de inteiro) $\to$ CWE-119 (violação dos
limites de \textit{buffer} de memória).
Em hardware Pixhawk sem MMU, uma escrita fora dos limites poderia
corromper a memória adjacente; sob o modelo de memória plano do
NuttX, ela não precisa gerar uma falha limpa, embora a
alcançabilidade no driver de produção, estabelecida a seguir, limite
o impacto físico.

Como o mantenedor do PX4 (mrpollo) mostrou na revisão do PR~\#27576,
e como a análise do código-fonte confirmou em seguida, o driver de
produção responde a essa questão negativamente: o laço de escrita
satura \texttt{write\_len} ao espaço livre restante; assim, o nível
de preenchimento nunca excede o \textit{buffer} antes que o registro
seja publicado e o campo seja zerado. A violação é, portanto,
alcançável no modelo, mas não no driver de produção. Nós a mantemos
como exemplo de por que um contraexemplo formalmente válido precisa
ser confrontado com o código de produção antes de ser chamado de
defeito.
\begin{rqbox}
\textbf{Esses resultados respondem afirmativamente à QP2} para pelo
menos um módulo: a BMC expôs uma violação de segurança de memória
alcançável no modelo, no \textit{harness} do driver de GPS, que os
testes unitários existentes do PX4 não exercitam.
\end{rqbox}

\begin{lstlisting}[float=t,basicstyle=\ttfamily\scriptsize,
numbers=none,language={},captionpos=b,
caption={Contraexemplo do ESBMC: \textit{underflow} sem sinal em
\texttt{dumpGpsData()} se propaga para um \texttt{memcpy} fora dos
limites (CWE-191 $\rightarrow$ CWE-119). O \texttt{dump\_buffer}
local do \textit{harness} é o \textit{buffer} \texttt{dump\_data}
passado ao driver.},label=lst:bug-gps]
[Counterexample]
State 3 -- test_gps_real_bit_operation, line 179
  input_len = 10
State 8 -- test_gps_real_bit_operation, line 185
  dump_buffer.len = 207
State 21 -- __memcpy_impl, line 280
  Violated property: dereference failure:
  Access to object out of bounds

VERIFICATION FAILED
\end{lstlisting}

\paragraph{Verificação da correção}
No \textit{harness} corrigido, o CBMC~5.95.1 reporta $0/7$
propriedades violadas (PR~\#27576; \textit{log} no artefato). Esta é
uma verificação da proteção proposta no \textit{harness}, e não uma
evidência sobre o driver de produção: o \textit{harness} do CBMC sem
a correção reporta os mesmos $0/7$ sob a mesma pré-condição de
comprimento limitado, de modo que a comparação não isola o efeito da
proteção. O ESBMC falha no \textit{harness} sem a correção
(Tabela~\ref{tab:results}); o artefato não tem \textit{log} do ESBMC
para o \textit{harness} corrigido.

\subsection{Discussão}
\label{subsec:discussion}

\paragraph{Configuração, correção e tempo}
\begin{rqbox}
\textbf{Esses resultados respondem à QP3}: três configurações por
módulo bastam; MAVLink com \texttt{--unwind~300} (laços profundos de
\textit{parsing}), VTOL com \texttt{--no-library} (desabilita a
biblioteca C abstrata) e Navigator com \texttt{--unwind~15}, todas
documentadas no repositório de artefatos.
\end{rqbox}
Um \texttt{grep} na árvore de testes do PX4 não encontra nenhum teste
unitário que referencie \texttt{dumpGpsData} ou a classe \texttt{GPS}
(isso indica apenas a ausência de cobertura \emph{direta}; testes de
\textit{software-in-the-loop} (SITL) ainda podem exercitá-la); assim,
a BMC e os testes unitários continuam
complementares~\cite{Godefroid2005,Khatiri2024}.

\paragraph{Observação em outro piloto automático e CI/CD}
A mesma observação de falta de validação de faixa (CWE-1284) aparece
no driver equivalente do ArduPilot~\cite{ArduPilot}: executar o
CBMC~5.95.1 sobre \texttt{AP\_InertialSensor\_BMI088.cpp}
(linhas~383--385) a reproduz (ArduPilot
Issue~\#33272\footnote{\url{https://github.com/ArduPilot/ardupilot/issues/33272}}).
Um mantenedor do ArduPilot encerrou a \textit{issue}, citando a faixa
de medição de $-104$ a $150\,^\circ$C do \textit{datasheet}. Os dois
projetos a recusaram; nós a reportamos apenas como uma verificação de
consistência.
Um \textit{workflow} do GitHub Actions em um \textit{fork} do
PX4\footnote{\url{https://github.com/marcela2881/PX4-Autopilot}}
executa os \textit{harnesses} da IMU e do GPS a cada \textit{push} e
\textit{pull request}.

\section{Ameaças à Validade}
\label{sec:threats}

\paragraph{Validade interna}
Os tempos reportados (medianas de três execuções) indicam apenas
ordens de grandeza.
A triagem descartou outros dois candidatos iniciais (MAVLink e EKF2)
como artefatos de configuração; seus \textit{logs} não estão no
artefato. O \textit{underflow} do GPS foi encontrado com o Boolector;
o Z3 não retornou veredito dentro do tempo limite (assim como no
MAVLink e no EKF2 com \texttt{--unwind~10} em nossas reexecuções com
a versão 8.1); portanto, o resultado depende do resolvedor; a
verificação cruzada com o CBMC é parcial para o GPS (limite
reduzido).

\paragraph{Validade externa}
Selecionar módulos pela viabilidade dos \textit{stubs} sub-representa
os módulos com forte acoplamento ao hardware ou ao escalonador; assim,
a amostra não é estatisticamente representativa do PX4. A análise
modular não detecta bugs de integração decorrentes de interações
concorrentes.

\paragraph{Validade de construto}
A verificação foi limitada a $k=10$ e, com
\texttt{--no-unwinding-assertions} habilitado, os laços que excedem o
limite são truncados silenciosamente; assim, um veredito
\textsc{successful} pode esconder falsos negativos mesmo abaixo da
profundidade operacional pretendida.

\paragraph{Correção dos \textit{stubs}}
Os \textit{stubs} que retornam \texttt{nondet\_T()} sobreaproximam as
\emph{entradas} do hardware; assim, um \textit{harness} que omite as
proteções de produção pode admitir estados que o driver real rejeita,
como mostrou o candidato da FIFO; a genuinidade depende da
confrontação com o código-fonte de produção. Por outro lado, os
\textit{stubs} \emph{subaproximam} o comportamento interno do
hardware (máquinas de estado de I2C/SPI, escalonamento do NuttX);
assim, defeitos disparados por esse estado interno não seriam
detectados.

\section{Conclusão e Trabalhos Futuros}
\label{sec:conclusion}

Apresentamos uma metodologia sistemática de verificação aplicada a
\textit{harnesses} de doze módulos do PX4 Autopilot (\textit{commit}
\texttt{8be2590c15}); seu resultado principal é uma garantia
limitada, em nível de modelo, com os \textit{harnesses} de dez
módulos, que abrangem classes representativas de entradas não
confiáveis, verificados livres de violações de segurança de memória
e de \textit{overflow} até o limite (${\approx}94$\,s no total;
${\approx}159$\,s com o GPS). O ESBMC também revelou uma violação de
segurança de memória no \textit{harness} do driver uBlox de GPS
(Issue~\#26866), reproduzida com o CBMC sob um \textit{harness}
semeado e com limite reduzido, e proposta ao projeto original com uma
verificação de limites (PR~\#27576); a revisão do mantenedor do PX4,
confirmada pela análise do código-fonte do driver de produção,
mostrou em seguida que o estado violador está fora de alcance nesse
driver e enquadrou a verificação como um reforço defensivo.
O candidato da FIFO do BMI088 foi um falso positivo na inspeção do
código-fonte, e a observação de temperatura foi recusada.
Trabalhos futuros: \textit{k}-indução, os módulos \textit{commander},
\textit{failsafe} e \textit{geofence}, e CI do PX4 com avaliação
\textit{Hardware-in-the-Loop} (HIL).

\section*{Agradecimentos}
M.A.R. agradece ao Prof. Lucas C. Cordeiro e ao Prof. Jos\'e
Reginaldo H. de Carvalho pela orientação. Os autores agradecem ao
projeto SAGRES e à comunidade de desenvolvimento do PX4. Este
trabalho foi realizado na Universidade Federal do Amazonas (UFAM),
com apoio de bolsa da CAPES (Código de Financiamento 001). IA
generativa (Claude, Anthropic) foi usada para análise de código e
para a redação e revisão do texto; os autores revisaram todo o
conteúdo e assumem a responsabilidade por ele.

\bibliographystyle{IEEEtran}
{\scriptsize\bibliography{references}}

\end{document}
